% Generated draft: requires mathematical and source review.
\documentclass[11pt]{article}
\usepackage[margin=25mm]{geometry}
\usepackage{fontspec}
\setmainfont{Latin Modern Roman}
\usepackage{amsmath,amssymb,mathtools,unicode-math}
\setmathfont{Latin Modern Math}
\usepackage{enumitem,needspace,xurl,hyperref}
\setlength{\parindent}{0pt}
\setlength{\parskip}{5pt}
\setlength{\emergencystretch}{3em}
\newcommand{\N}{\mathbb{N}}
\newcommand{\Z}{\mathbb{Z}}
\newcommand{\Q}{\mathbb{Q}}
\newcommand{\R}{\mathbb{R}}
\newcommand{\C}{\mathbb{C}}
\newcommand{\F}{\mathbb{F}}
\newcommand{\A}{\mathbb{A}}
\newcommand{\PP}{\mathbb{P}}
\newcommand{\E}{\mathbb{E}}
\newcommand{\eps}{\varepsilon}
\newcommand{\dd}{\,\mathrm{d}}
\begin{document}
\section*{UnsolvedMath ID: 30005041}
\textbf{Source problem code:} catalogue entry
\subsection{Definitions and mathematical statement}
Let \(\sigma\) be a fixed continuous homomorphism as in the supplied statement,
\[
\sigma:G\longrightarrow \operatorname{Aut}(X,[\mu])\ltimes L^0(X,\mu;\{-1,1\}).
\]
Here \([\mu]\) denotes the measure class of \(\mu\); under the usual interpretation, \(\operatorname{Aut}(X,[\mu])\) consists of invertible measurable transformations preserving this measure class, identified modulo equality almost everywhere, and \(L^0(X,\mu;\{-1,1\})\) consists of measurable sign functions identified modulo equality almost everywhere. For each finite real exponent \(p>0\), let \(\pi_{p,\sigma}\) denote the induced linear isometric representation on \(L^p(X,\mu)\) specified by the source. The same homomorphism \(\sigma\) is held fixed as \(p\) varies.

To make the cohomological target explicit, under the continuous, unreduced interpretation of first cohomology, put \(E_p=L^p(X,\mu)\) and define
\[
\begin{aligned}
Z^1(G,\pi_{p,\sigma})
&=\{b:G\to E_p:\ b\text{ is continuous and }\\[-2pt]
&\hspace{35mm}b(gh)=b(g)+\pi_{p,\sigma}(g)b(h)\text{ for all }g,h\in G\},\\
B^1(G,\pi_{p,\sigma})
&=\{b_v:\ b_v(g)=\pi_{p,\sigma}(g)v-v,\ v\in E_p\},\\
H^1(G,\pi_{p,\sigma})
&=Z^1(G,\pi_{p,\sigma})/B^1(G,\pi_{p,\sigma}).
\end{aligned}
\]
This definition presupposes that the induced representations have the continuity needed for these coboundaries to be continuous. Vanishing means that every continuous cocycle is a coboundary; no closure of \(B^1\) is taken.

The question is whether, for every datum covered by the source, the set
\[
S_\sigma=\{p\in(0,\infty):H^1(G,\pi_{p,\sigma})=0\}
\]
is an interval. Explicitly, is it always true that
\[
\begin{gathered}
0<p_0<q<p_1<\infty,\\
H^1(G,\pi_{p_0,\sigma})=0,
\qquad H^1(G,\pi_{p_1,\sigma})=0
\quad\Longrightarrow\quad
H^1(G,\pi_{q,\sigma})=0?
\end{gathered}
\]
This order-convex formulation permits the empty set and singletons and imposes no endpoint inclusion condition.

The supplied record does not specify the class of topological groups \(G\), hypotheses on \((X,\mu)\), the scalar field of \(L^p\), the topology on the target of \(\sigma\), or the precise construction and continuity assumptions for the induced representations. It also does not explicitly specify the cohomology convention; the continuous, unreduced convention above therefore requires confirmation. The range \(0<p<1\) is retained: these spaces generally have a quasi-norm rather than a norm, so a Banach-space-only interpretation would not cover the stated question. These unresolved qualifications must be checked before this becomes a fully specified statement.
\subsection{Short English statement}
For one fixed homomorphism \(\sigma\), must vanishing of first cohomology at two positive exponents force vanishing at every exponent between them? The supplied record presents this as a refined form of Drutu's conjecture and labels it open. Its analytic hypotheses and cohomology convention require confirmation.
\subsection{Sources}
\begin{itemize}
\item Supplied local UnsolvedMath registry record, permanent ID \(30005041\), title \emph{Interval Structure of Cohomology-Vanishing Exponents}, exact source problem code \texttt{catalogue entry}, in \emph{Oberwolfach Reports Open Problems}. No individual proposer is supplied; the proposed year is recorded as 2022.
\item Supplied report citation: \emph{Geometric Structures in Group Theory (2023)}, Oberwolfach Reports, DOI \texttt{10.4171/owr/2022/11}, pp.\ 517--576. \url{https://doi.org/10.4171/owr/2022/11}. The background describes these reports as documenting mathematical workshops and their open-problem discussions. The supplied contextual passage calls this question a refined form of Drutu's conjecture; it does not supply that conjecture's full formulation.
\item Supplied literature reference: \emph{Isometric actions on Lp spaces}. \url{https://arxiv.org/abs/2001.02490}.
\item Supplied literature reference: \emph{Oberwolfach report}. \url{https://doi.org/10.4171/owr/2022/11}.
\item Literature assessment dated exactly \texttt{2026-08-22}: the supplied curation labels the problem open and states that the cited work explicitly leaves the general fixed-representation interval question unsettled. This is supplied dated curation, not an independently verified literature assessment; the background expressly requests independent verification of time-sensitive status and cited claims.
\item The registry labels the statement \texttt{corrected\_verified} and records that the official report was checked directly and the statement restored while removing extraction artifacts. These are supplied verification claims, not verification performed for this draft.
\end{itemize}
\end{document}
